\section{Provenance and Study Status}
\subsection{Frozen final iteration and information boundaries}
Exactly three new scientific analyses were permitted: all ordered cross-architecture phase contrasts, acquired-only neural influence, and a fixed-context perturbation test of Proposition 1. Original papers and inputs were preserved before new inference. This is a frozen retrospective iteration on reused cohorts, not preregistration. The optional unused-cohort branch is skipped because no authenticated untouched cohort was established. Administrative submission metadata are outside this scientific iteration.

The SAINT coordinate amendment was recorded after interface inspection but before any new inference: both new phases use physical $z/(\mathrm{source\ depth}-1)$. Legacy coordinates are preserved for reproduction. New FP32 finite-response probes also differ from historical mixed precision. These are explicit implementation contracts, not favorable-result substitutions. Raw arrays and the initial/amended protocol hashes are retained.

\subsection{Provenance, Implementation Correction, and Claim Levels}
This revision is a retrospective synthesis, not a new preregistered confirmatory study. The machine-readable research ledger records source, split, checkpoint and training-configuration provenance. Candidate availability does not establish checkpoint usage: the ten main five-architecture checkpoints are separately bound to inference code and SHA-256 hashes. Recorded manifests, source lines and unsuccessful analyses are preserved. The only recovered paper Git commit includes both historical architecture-family protocols and results; it does not establish a pre-result commit sequence. Filesystem timestamps and protocol assertions cannot close that gap.

The original historical architecture-family ranking-primary hypothesis preceded the later narrative emphasis on local observability according to the documents, but a complete independently verifiable chronology is unavailable. The reversal-control amendment followed an unfavorable one-case smoke outcome. Its construct-validity rationale is recorded; it is not retroactively called outcome-independent. external MRA bright-anchor support was motivated after inspecting TopCoW results; CAS was a subsequent fixed-rule transfer to another public challenge training collection, not its hidden test set. All analyses now using these cohorts are aware of their historical results.

The fastMRI split contains 112/24/24 training/validation/test series and 111/24/24 corresponding patient IDs. One patient overlaps training and validation; test IDs are disjoint. IXI has 404/58/115 disjoint IDs. Eight NingXia patient folds have no within-fold train/validation/test patient overlap. These checks verify local identifiers, not cross-dataset patient linkage.

\paragraph{A corrected influence definition.}
Earlier manuscript text described averaging response cosines. Both the main HAT/SwinIR implementation and the historical architecture-family new-architecture implementation actually sum shared-anchor response inner products and normalize by \emph{all}-anchor response energy, as in the revised main equation. Historical values remain unchanged. For example, responses $(1,0)$ on a shared anchor and $(0,3)$ versus $(0,4)$ on two different unshared anchors give embedded cosine $1/\sqrt{170}$, whereas the shared-only response cosine is one. Omitting the unshared response energy would define a materially different statistic.

The historical original clean HAT/SwinIR correlation pools plane/model observations against a fixed endpoint (phase-0 MAE for fastMRI/NingXia; disjoint-pair mean MAE for IXI) while resampling examination/patient clusters. It must not be described as an examination-level point correlation. The five-architecture ranking analysis instead first aggregates within examinations. Those estimands remain separate in the revised figures.

\paragraph{Reference self-pair sensitivity.}
Final source review also found that the historical broad-grid summaries include the reference-versus-itself context, for which overlap count is two, influence is one, and error cosine is one. We retain those original tables below but add an explicitly post-hoc integrity sensitivity excluding all 966 such rows. The remaining 7,728 rows yield influence-minus-count alignment advantages .061/.085 on fastMRI, .128/.118 on IXI, and .070/.060 on NingXia (HAT/SwinIR). All six-stratum simultaneous intervals are positive. The bootstrap uses 5,000 examination/patient draws, shared across models within a cohort and independent streams across cohorts. Weighted-rank calculations are verified against direct repeated-row resampling. The preserved historical non-self figure and table below use this result; the final main Figure 3 and Table 2 instead report the new acquired-only experiment. Historical fixed-one-overlap estimates are unchanged. This strengthens a scoped beyond-count association, not a prospective or causal claim.

\paragraph{Claim hierarchy.}
Level A is replicated local error-observability loss in the evaluated designs. Level B comprises influence beyond count where stable, case-ranking effects, and the matched-geometry interpretation. Level C comprises universal scalar superiority, prospective acquisition benefit, universal target-free neural mechanism validation, and non-MRI generalization. Level A denotes reproducibility in these data, not prespecification or clinical validation. Exact claim-to-experiment, unit, inference and figure mappings are supplied in the final reproducibility and evidence map.

\subsection{Scope and Evidence Contract}
This supplement reports protocol details, full strata, and every result that materially limits a submitted claim. Experiments are included according to relevance, not favorability. The manuscript does not claim a new reconstruction architecture, calibrated uncertainty, universal estimator superiority, plane-invariant influence, local failure localization, or prospective acquisition benefit.

The independent units are 24 fastMRI examinations, 115 IXI examinations, and eight NingXia patients containing 22 series. All NingXia intervals resample patients with subordinate series, planes, models, and contexts retained. Missing target values are not subtracted in a disagreement score, but its evaluation support can be target-derived. The historical architecture studies use a ground-truth 1\% mask; the recovered bank uses the target model's prediction mask, not ground truth. Anchor-CV uses held-out acquired anchors (including their masks); adjacent MAE uses masks from acquired pairs. These information regimes must not be conflated.


\subsection{Data Attribution and Unresolved Author Documentation}
The brain T2 collection is locally catalogued as NYU fastMRI and converted from local DICOM series to NIfTI. The conversion manifest and physical-position/pixel-order checks trace the analyzed files, but do not authenticate an official download release or accession. We retain the fastMRI initiative citation required by its data portal, and add the brain challenge publication for brain-dataset context; we do not claim that our local split is the official challenge split. The author must confirm the precise release, access agreement and image-publication permission before submission. NYU fastMRI investigators supplied the attributed resource, without participating in this analysis or manuscript; this attribution is conditional on confirmation of local provenance.

IXI is attributed to the Biomedical Image Analysis Group at Imperial College London. Its public T2 NIfTI collection is described on the Brain Development website and distributed under CC BY-SA 3.0. Our local training/validation/test partition is separate from the source collection. No unsupported dataset publication year is assigned.

NingXia denotes a locally retained knee T2 collection with eight patient groups and 22 analyzed series. Acquisition institution, data custodian, release/allowed-use agreement, de-identification process, IRB approval or exemption, and consent or consent-waiver status are not established by the repository documents reviewed. Local patient-group splitting establishes analysis clustering, not ethics authorization or de-identification. \textbf{Author confirmation required before submission:} supply those records and the verified statement here; no approval number, consent waiver or permission is inferred. Direct identifiers found in raw directory naming are not reproduced in the manuscript or review package, and de-identification of the research extract cannot presently be certified.

TopCoW and CAS remain public-challenge training-collection analyses with locally available labels, not hidden-test evaluations. Dataset citations and the remaining local permission/release checks are listed in the final evidence map. No new subject data were acquired for this submission patch.
